\chapter{Compréhension des Données - Data Understanding}

\begin{objectivebox}
Cette phase CRISP-DM se concentre sur la collecte initiale des données, leur exploration et leur qualification pour comprendre la nature et la qualité des informations disponibles pour le projet Housy Tunisia.
\end{objectivebox}

\section{Collecte des Données Initiales}

\subsection{Sources de Données Identifiées}

Pour le projet Housy Tunisia, plusieurs sources de données ont été identifiées et exploitées :

\begin{enumerate}
    \item \textbf{Données de projets de construction :}
    \begin{itemize}
        \item Historique des projets réalisés
        \item Types de construction (villas, immeubles, bureaux)
        \item Surfaces et caractéristiques techniques
        \item Coûts réalisés et délais
    \end{itemize}
    
    \item \textbf{Données des matériaux de construction :}
    \begin{itemize}
        \item Catalogues fournisseurs tunisiens
        \item Prix des matériaux par région
        \item Évolution historique des prix
        \item Spécifications techniques
    \end{itemize}
    
    \item \textbf{Données de main-d'œuvre :}
    \begin{itemize}
        \item Tarifs par corps de métier
        \item Productivité moyenne par région
        \item Disponibilité des compétences
        \item Coûts sociaux et charges
    \end{itemize}
    
    \item \textbf{Données réglementaires :}
    \begin{itemize}
        \item Normes de construction tunisiennes
        \item Réglementations locales par gouvernorat
        \item Procédures administratives
        \item Taxes et frais applicables
    \end{itemize}
\end{enumerate}

\textbf{[CAPTURE À AJOUTER : Schéma des sources de données et leurs interconnexions]}

\section{Structure de la Base de Données}

La base de données Housy Tunisia est organisée en 43 tables réparties en 6 domaines fonctionnels :

\subsection{Domaine Gestion de Projets (10 tables)}

\begin{itemize}
    \item \texttt{users} : Utilisateurs du système
    \item \texttt{projects} : Projets de construction
    \item \texttt{tasks} : Tâches et activités
    \item \texttt{resources} : Ressources allouées
    \item \texttt{project\_phases} : Phases de réalisation
    \item \texttt{project\_estimations} : Estimations de coûts
\end{itemize}

\textbf{[CAPTURE À AJOUTER : Diagramme ERD du domaine Gestion de Projets]}

\subsection{Volume et Couverture des Données}

\begin{table}[h]
\centering
\begin{tabular}{|l|c|c|l|}
\hline
\textbf{Type de Données} & \textbf{Volume} & \textbf{Période} & \textbf{Source} \\
\hline
Projets de construction & 1,200+ entrées & 2020-2025 & Bases internes + partenaires \\
\hline
Matériaux de construction & 800+ références & 2023-2025 & Catalogues fournisseurs \\
\hline
Prix main-d'œuvre & 50+ métiers & 2024-2025 & Enquêtes terrain \\
\hline
Données géographiques & 24 gouvernorats & Permanent & Données officielles \\
\hline
\end{tabular}
\caption{Volume et couverture des données collectées}
\end{table}

\section{Exploration des Données}

\subsection{Répartition des Projets par Type}

\begin{table}[h]
\centering
\begin{tabular}{|l|c|c|c|}
\hline
\textbf{Type de Projet} & \textbf{Nombre} & \textbf{Pourcentage} & \textbf{Coût Moyen (TND)} \\
\hline
Villa Individuelle & 450 & 37.5\% & 125,000 \\
\hline
Immeuble Résidentiel & 320 & 26.7\% & 550,000 \\
\hline
Bureau Commercial & 280 & 23.3\% & 210,000 \\
\hline
Infrastructure & 150 & 12.5\% & 1,200,000 \\
\hline
\textbf{Total} & \textbf{1,200} & \textbf{100\%} & \textbf{346,750} \\
\hline
\end{tabular}
\caption{Distribution des projets par type dans la base de données}
\end{table}

\textbf{[CAPTURE À AJOUTER : Graphique en secteurs de la répartition des projets]}

\subsection{Analyse de la Qualité des Données}

\begin{table}[h]
\centering
\begin{tabular}{|l|c|c|c|}
\hline
\textbf{Champ} & \textbf{Complétude} & \textbf{Valeurs Manquantes} & \textbf{Qualité} \\
\hline
Nom du projet & 100\% & 0 & Excellente \\
\hline
Surface totale & 98.5\% & 18 & Très bonne \\
\hline
Coût total & 96.2\% & 46 & Bonne \\
\hline
Date de fin & 92.1\% & 95 & Acceptable \\
\hline
\end{tabular}
\caption{Analyse de complétude des données}
\end{table}

\textbf{[CAPTURE À AJOUTER : Dashboard de qualité des données]}

\section{Types de Construction Tunisiens}

\subsection{Spécificités Locales}

\begin{itemize}
    \item \textbf{Villa traditionnelle :} Architecture arabo-andalouse
    \item \textbf{Villa moderne :} Design contemporain
    \item \textbf{Immeuble R+4 :} Standard urbain tunisien
    \item \textbf{Logement social :} Programmes gouvernementaux
\end{itemize}

\textbf{[CAPTURE À AJOUTER : Galerie des types de construction tunisiens]}

\begin{resultbox}
Cette phase d'exploration des données a permis de comprendre la structure, la qualité et les caractéristiques des données disponibles. La base constituée de 43 tables et plus de 1,200 projets de référence offre une foundation solide pour les phases suivantes.
\end{resultbox}
